\documentclass[aspectratio=169,10pt]{beamer}
\usetheme{skd}

\title[Pertemuan 2 -- Pencacahan]{Sistem Kontrol Digital: Pencacahan dan Rekonstruksi}
\subtitle{Digital Control Systems: Sampling and Reconstruction}
\author{Ir. Estiyanti Ekawati, M.T., Ph.D. \quad\textbar\quad Faqihza Mukhlish, S.T., M.T., Ph.D.}
\institute{Program Studi Teknik Fisika, Fakultas Teknologi Industri, Institut Teknologi Bandung}
\date{Pertemuan 2 \textbar\ Semester I 2026/2027}

\begin{document}

\begin{frame}[plain,noframenumbering]
	\titlepage
\end{frame}

\begin{frame}{Tujuan pertemuan}
	\framesubtitle{Learning objectives}
	\begin{columns}[T]
		\begin{column}{0.55\textwidth}
			Setelah pertemuan ini Anda mampu:
			\begin{enumerate}
				\item membedakan sinyal kontinu, diskrit, dan digital;
				\item menghitung resolusi dan derau kuantisasi ADC/DAC;
				\item menjelaskan spektrum sinyal tercacah dan teorema Shannon;
				\item mengenali dan mencegah \textit{aliasing};
				\item menurunkan karakteristik \textit{zero-order hold};
				\item memilih waktu cacah $T$ secara beralasan.
			\end{enumerate}
		\end{column}
		\begin{column}{0.4\textwidth}
			\begin{block}{Sub-CPMK S2}
				\small Menjelaskan proses pencacahan, kuantisasi, rekonstruksi, dan \textit{aliasing}, serta memilih waktu cacah.\\[3pt]
				\en{Explain sampling, quantisation, reconstruction, and aliasing, and select a sampling period.}
			\end{block}
			{\small Bacaan: [FR] 8.1, 8.4, 8.5; [OG] 1.4--1.5, 3.2; [PH] 3}
		\end{column}
	\end{columns}
\end{frame}

% ===========================================================================
\section{Rantai Konversi Sinyal\texorpdfstring{\\[4pt]{\Large\en{Signal Conversion Chain}}}{ / Signal Conversion Chain}}
% ===========================================================================

\begin{frame}{Dari dunia fisik ke komputer dan kembali}
	\framesubtitle{From the physical world to the computer and back}
	\centering
	\begin{tikzpicture}[node distance=0.35cm and 0.45cm]
		\node[blk] (sen) {Sensor};
		\node[blk,right=of sen] (aa) {Filter\\\textit{anti-alias}};
		\node[blk,right=of aa] (sh) {\textit{Sample}\\\& \textit{Hold}};
		\node[blk,right=of sh] (adc) {ADC\\(kuantisasi)};
		\node[blkd,right=of adc] (cpu) {Algoritma\\kontrol};
		\node[blk,right=of cpu] (dac) {DAC\\+ ZOH};
		\node[blk,right=of dac] (act) {Aktuator};
		\draw[sig] (sen) -- (aa); \draw[sig] (aa) -- (sh); \draw[sig] (sh) -- (adc);
		\draw[sig] (adc) -- (cpu); \draw[sig] (cpu) -- (dac); \draw[sig] (dac) -- (act);
		\node[lbl,below=0.15cm of aa] {kontinu};
		\node[lbl,below=0.15cm of sh] {diskrit};
		\node[lbl,below=0.15cm of cpu] {digital};
		\node[lbl,below=0.15cm of dac] {kontinu bertangga};
	\end{tikzpicture}

	\vspace{0.5cm}
	\begin{columns}[T]
		\begin{column}{0.48\textwidth}
			\begin{itemize}
				\item \term{Pencacahan}{sampling}: mengambil nilai pada $t=kT$ (kuantisasi waktu).
				\item \term{Kuantisasi}{quantisation}: membulatkan amplitudo ke tingkat terdekat.
			\end{itemize}
		\end{column}
		\begin{column}{0.48\textwidth}
			\begin{itemize}
				\item \term{Rekonstruksi}{reconstruction}: mengubah barisan angka kembali menjadi sinyal kontinu, umumnya dengan ZOH.
				\item Kit SKD: ADC 10-bit internal, DAC MCP4725 12-bit.
			\end{itemize}
		\end{column}
	\end{columns}
\end{frame}

\begin{frame}{Kuantisasi amplitudo}
	\framesubtitle{Amplitude quantisation}
	\begin{columns}[c]
		\begin{column}{0.48\textwidth}
			\centering
			\begin{tikzpicture}
				\begin{axis}[skd,width=7cm,height=5.4cm,xmin=0,xmax=1,ymin=0,ymax=1,xlabel={masukan analog / FSR},ylabel={keluaran / FSR},xtick={0,0.25,0.5,0.75,1},ytick={0,0.25,0.5,0.75,1}]
					\addplot[skdgray,dashed,domain=0:1] {x};
					\addplot[skdblue,const plot,samples at={0,0.0625,...,1}] {floor(x*8+0.5)/8};
				\end{axis}
			\end{tikzpicture}\\
			{\scriptsize Kuantiser 3-bit (8 tingkat), pembulatan}
		\end{column}
		\begin{column}{0.48\textwidth}
			\small
			Untuk ADC $n$-bit dengan rentang penuh FSR:
			\[
				Q = \text{LSB} = \frac{\text{FSR}}{2^n}
			\]
			Galat kuantisasi (pembulatan): $|e_q| \le Q/2$. Bila sinyal cukup ``sibuk'', $e_q$ dapat dimodelkan derau seragam:
			\[
				\sigma_q^2 = \frac{1}{Q}\int_{-Q/2}^{Q/2} \xi^2\,d\xi = \frac{Q^2}{12}
			\]
			\scriptsize
			\renewcommand{\arraystretch}{1.2}
			\begin{tabular}{@{}l c c@{}}
				\toprule
				Perangkat & $Q$ & $\sigma_q$ \\
				\midrule
				ADC Arduino 10-bit, 5\,V & 4,88\,mV & 1,41\,mV \\
				DAC MCP4725 12-bit, 5\,V & 1,22\,mV & 0,35\,mV \\
				ADC 16-bit, 10\,V & 0,15\,mV & 0,04\,mV \\
				\bottomrule
			\end{tabular}
		\end{column}
	\end{columns}
\end{frame}

\begin{frame}{Dampak kuantisasi pada kontrol}
	\framesubtitle{Effects of quantisation on control}
	\begin{columns}[T]
		\begin{column}{0.48\textwidth}
			\begin{itemize}
				\item \textbf{Batas resolusi}: perubahan di bawah $Q$ tidak terlihat; galat tunak minimum $\approx Q/2$.
				\item \textbf{Siklus batas}: kontroler dengan integrator dapat ``berburu'' bolak-balik antar dua tingkat kuantisasi.
				\item \textbf{Derivatif}: $\Delta y/T$ hanya bernilai kelipatan $Q/T$; makin kecil $T$, makin kasar estimasi turunan.
			\end{itemize}
		\end{column}
		\begin{column}{0.48\textwidth}
			\begin{exampleblock}{Contoh kit}
				\small
				Potensiometer satu putaran penuh $=5$\,V, ADC 10-bit $\Rightarrow$ resolusi $360^\circ/1024 \approx 0{,}35^\circ$ poros potensiometer.\\[3pt]
				Dengan $T=10$\,ms, satu LSB setara kecepatan semu $0{,}35^\circ/0{,}01\,\text{s} = 35^\circ$/s.
			\end{exampleblock}
			{\small Solusi umum: perataan beberapa sampel (fungsi \texttt{smooth\_...} di program praktikum), filter, atau estimator keadaan (Pertemuan 13).}
		\end{column}
	\end{columns}
\end{frame}

% ===========================================================================
\section{Pencacahan dan Teorema Shannon\texorpdfstring{\\[4pt]{\Large\en{Sampling and Shannon's Theorem}}}{ / Sampling and Shannon's Theorem}}
% ===========================================================================

\begin{frame}{Model matematis pencacahan}
	\framesubtitle{Mathematical model of sampling}
	\begin{columns}[c]
		\begin{column}{0.5\textwidth}
			\small
			Pencacah ideal = perkalian dengan deret impuls:
			\[
				x^*(t) = x(t)\sum_{k=-\infty}^{\infty}\delta(t-kT) = \sum_k x(kT)\,\delta(t-kT)
			\]
			Transformasi Laplace:
			\[
				X^*(s) = \sum_{k=0}^{\infty} x(kT)\,e^{-kTs}
			\]
			Dengan substitusi $z = e^{sT}$ diperoleh transformasi Z, $X(z)=\sum x(kT)z^{-k}$ (Pertemuan 3).
		\end{column}
		\begin{column}{0.46\textwidth}
			\centering
			\begin{tikzpicture}
				\begin{axis}[skd,width=6.8cm,height=4.4cm,xmin=0,xmax=6,ymin=-1.3,ymax=1.3,xlabel={$t$},ytick=\empty,xtick={0,1,2,3,4,5,6},xticklabels={0,$T$,$2T$,$3T$,$4T$,$5T$,$6T$}]
					\addplot[skdblue!40,domain=0:6,samples=100] {sin(deg(1.1*x))};
					\addplot[skdorange,ycomb,mark=triangle*,mark size=2pt,samples at={0,1,...,6}] {sin(deg(1.1*x))};
				\end{axis}
			\end{tikzpicture}\\
			{\scriptsize $x^*(t)$: impuls dengan bobot $x(kT)$}
		\end{column}
	\end{columns}
\end{frame}

\begin{frame}{Spektrum sinyal tercacah}
	\framesubtitle{Spectrum of a sampled signal}
	\begin{columns}[c]
		\begin{column}{0.55\textwidth}
			\centering
			\begin{tikzpicture}[yscale=0.9]
				\foreach \y/\ws/\c in {0/3.2/skdblue, -2.3/1.6/skdred} {
						\draw[-{Stealth}] (-4.3,\y) -- (4.3,\y) node[right,font=\scriptsize]{$\omega$};
						\foreach \m in {-2,-1,0,1,2} {
								\pgfmathsetmacro{\cx}{\m*\ws}
								\ifdim \cx pt > -4.2pt \ifdim \cx pt < 4.2pt
									\fill[\c!25,draw=\c] (\cx-1.0,\y) -- (\cx,\y+1.2) -- (\cx+1.0,\y) -- cycle;
								\fi\fi
							}
						\draw (0,\y) -- (0,\y-0.1) node[below,font=\tiny]{0};
						\draw (\ws,\y) -- (\ws,\y-0.1) node[below,font=\tiny]{$\omega_s$};
						\draw (\ws/2,\y) -- (\ws/2,\y-0.1) node[below,font=\tiny]{$\omega_N$};
					}
				\node[font=\scriptsize,anchor=west] at (-4.3,1.25) {$\omega_s > 2\omega_m$: replika terpisah};
				\node[font=\scriptsize,anchor=west,text=skdred] at (-4.3,-1.05) {$\omega_s < 2\omega_m$: replika tumpang tindih (\textit{aliasing})};
			\end{tikzpicture}
		\end{column}
		\begin{column}{0.42\textwidth}
			\small
			Pencacahan membuat spektrum menjadi \textbf{periodik}:
			\[
				X^*(j\omega) = \frac{1}{T}\sum_{n=-\infty}^{\infty} X\big(j(\omega - n\omega_s)\big)
			\]
			dengan $\omega_s = 2\pi/T$.\\[4pt]
			\term{Frekuensi Nyquist}{Nyquist frequency}: $\omega_N = \omega_s/2$.
		\end{column}
	\end{columns}
\end{frame}

\begin{frame}{Teorema pencacahan Shannon}
	\framesubtitle{Shannon's sampling theorem}
	\begin{block}{Teorema}
		Sinyal $x(t)$ yang spektrumnya terbatas pada $|\omega| < \omega_m$ dapat direkonstruksi sempurna dari sampel $x(kT)$ jika
		\[
			\omega_s > 2\,\omega_m .
		\]
		\en{A band-limited signal is fully recoverable from its samples if the sampling frequency exceeds twice its highest frequency.}
	\end{block}
	\vspace{0.2cm}
	Rekonstruksi ideal (interpolasi sinc):
	\[
		x(t) = \sum_{k=-\infty}^{\infty} x(kT)\,\frac{\sin\big(\omega_s(t-kT)/2\big)}{\omega_s(t-kT)/2}
	\]
	\kunci{Rekonstruksi ini \alert{nonkausal} (memerlukan sampel masa depan), sehingga tidak dapat dipakai dalam kontrol waktu nyata. Kontroler memakai \textbf{ZOH}.}
\end{frame}

\begin{frame}{\textit{Aliasing}: frekuensi yang menyamar}
	\framesubtitle{Aliasing: frequencies in disguise}
	\begin{columns}[c]
		\begin{column}{0.58\textwidth}
			\centering
			\begin{tikzpicture}
				\begin{axis}[skd,width=9cm,height=5cm,xmin=0,xmax=1,ymin=-1.3,ymax=1.3,xlabel={$t$ (s)},ytick={-1,0,1},legend style={at={(0.5,1.02)},anchor=south},legend columns=3]
					\addplot[skdblue!60,domain=0:1,samples=400,thin] {sin(deg(2*pi*9*x))};
					\addplot[skdorange,dashed,domain=0:1,samples=100] {-sin(deg(2*pi*x))};
					\addplot[skdred,only marks,mark=*,mark size=1.6pt,samples at={0,0.1,...,1}] {sin(deg(2*pi*9*x))};
					\legend{sinyal 9 Hz, alias 1 Hz, sampel $f_s=10$ Hz}
				\end{axis}
			\end{tikzpicture}
		\end{column}
		\begin{column}{0.38\textwidth}
			\small
			Sinyal 9 Hz dicacah 10 Hz menghasilkan sampel yang \textbf{identik} dengan sinyal 1 Hz.\\[4pt]
			Frekuensi alias:
			\[
				f_a = \left| f - n f_s \right|,\quad n = \text{round}(f/f_s)
			\]
			\scriptsize
			\begin{tabular}{@{}c c c@{}}
				\toprule
				$f$ & $f_s$ & $f_a$ \\
				\midrule
				9 Hz & 10 Hz & 1 Hz \\
				60 Hz & 50 Hz & 10 Hz \\
				130 Hz & 100 Hz & 30 Hz \\
				\bottomrule
			\end{tabular}
		\end{column}
	\end{columns}
\end{frame}

\begin{frame}{Mencegah \textit{aliasing}}
	\framesubtitle{Preventing aliasing}
	\begin{columns}[T]
		\begin{column}{0.5\textwidth}
			\begin{itemize}
				\item Pasang \term{filter anti-alias}{anti-aliasing filter} (lolos rendah analog) \textbf{sebelum} ADC, dengan frekuensi potong di bawah $\omega_N$.
				\item Derau listrik 50 Hz dan derau komutasi motor dapat ber-alias menjadi frekuensi rendah yang mengganggu kalang.
				\item Filter anti-alias menambah tunda fasa; harus diperhitungkan dalam model plant.
				\item Alternatif: cacah sangat cepat ($T_{\text{cepat}}$), filter digital, lalu \textit{decimate} ke $T$ kontrol.
			\end{itemize}
		\end{column}
		\begin{column}{0.46\textwidth}
			\begin{exampleblock}{Pada kit SKD}
				\small
				Pengkondisi sinyal tachometer adalah penyearah + filter RC dengan $\tau_f = R_3C_4 = 0{,}1$\,s. Filter ini meredam riak dan derau tinggi, tetapi juga \alert{menambah dinamika} pada pengukuran kecepatan.\\[3pt]
				Pertanyaan: apa konsekuensinya bila waktu cacah dipilih jauh lebih kecil dari 0,1\,s?
			\end{exampleblock}
		\end{column}
	\end{columns}
\end{frame}

% ===========================================================================
\section{Rekonstruksi: \textit{Data Hold}\texorpdfstring{\\[4pt]{\Large\en{Reconstruction: Data Hold}}}{ / Reconstruction: Data Hold}}
% ===========================================================================

\begin{frame}{\textit{Zero-order hold} (ZOH)}
	\framesubtitle{Zero-order hold}
	\begin{columns}[c]
		\begin{column}{0.45\textwidth}
			\centering
			\begin{tikzpicture}
				\begin{axis}[skd,width=6.6cm,height=4.4cm,xmin=0,xmax=6,ymin=-1.3,ymax=1.3,xlabel={$t/T$},ytick=\empty]
					\addplot[skdblue!40,domain=0:6,samples=100] {sin(deg(1.1*x))};
					\addplot[skdgreen,const plot,samples at={0,1,...,6}] {sin(deg(1.1*x))};
					\addplot[skdorange,only marks,mark=*,mark size=1.3pt,samples at={0,1,...,6}] {sin(deg(1.1*x))};
				\end{axis}
			\end{tikzpicture}\\
			{\scriptsize $u(t) = u(kT)$ untuk $kT \le t < (k+1)T$}
		\end{column}
		\begin{column}{0.52\textwidth}
			\small
			Respons impuls ZOH: pulsa persegi selebar $T$,
			\[
				g_h(t) = 1(t) - 1(t-T)\ \Rightarrow\ G_h(s) = \frac{1 - e^{-sT}}{s}
			\]
			Respons frekuensi:
			\[
				G_h(j\omega) = T\,\frac{\sin(\omega T/2)}{\omega T/2}\,e^{-j\omega T/2}
			\]
			\begin{itemize}
				\item Penguatan $\approx T$ pada frekuensi rendah (diserap oleh pencacah $1/T$).
				\item Fasa $-\omega T/2$: ZOH \alert{setara tunda setengah waktu cacah}.
			\end{itemize}
		\end{column}
	\end{columns}
\end{frame}

\begin{frame}{Respons frekuensi ZOH}
	\framesubtitle{ZOH frequency response}
	\begin{columns}[c]
		\begin{column}{0.6\textwidth}
			\centering
			\begin{tikzpicture}
				\begin{groupplot}[group style={group size=1 by 2,vertical sep=0.5cm,xticklabels at=edge bottom},skd,width=8.6cm,height=3.4cm,xmin=0.01,xmax=3,xtick={0.5,1,1.5,2,2.5,3},xticklabels={$\tfrac{1}{2}\omega_s$,$\omega_s$,,$2\omega_s$,,$3\omega_s$}]
					\nextgroupplot[ylabel={$|G_h|/T$},ymin=0,ymax=1.1]
					\addplot[skdblue,domain=0.01:3,samples=300] {abs(sin(deg(pi*x))/(pi*x))};
					\draw[skdgray,dashed] (axis cs:0.5,0) -- (axis cs:0.5,1.1);
					\nextgroupplot[ylabel={fasa (\textdegree)},ymin=-560,ymax=10,xlabel={$\omega$},ytick={0,-180,-360,-540}]
					\addplot[skdorange,domain=0.01:3,samples=300] {-180*x + 180*floor(x)};
					\addplot[skdgray,dashed,domain=0.01:3] {-180*x};
				\end{groupplot}
			\end{tikzpicture}
		\end{column}
		\begin{column}{0.36\textwidth}
			\small
			\begin{itemize}
				\item Di $\omega_N$, magnitudo turun ke $2/\pi \approx 0{,}64$ ($-3{,}9$ dB).
				\item ZOH meloloskan sebagian replika spektral (bukan filter ideal).
				\item Garis putus: fasa $-\omega T/2$; lompatan $180^\circ$ terjadi di nol magnitudo.
			\end{itemize}
			\kunci{Pendekatan praktis: $G_h(s)/T \approx e^{-sT/2} \approx \dfrac{1}{1+sT/2}$.}
		\end{column}
	\end{columns}
\end{frame}

\begin{frame}{\textit{First-order hold} dan alternatif}
	\framesubtitle{First-order hold and alternatives}
	\begin{columns}[T]
		\begin{column}{0.5\textwidth}
			\small
			\textbf{FOH} mengekstrapolasi garis dari dua sampel terakhir:
			\[
				u(t) = u(kT) + \frac{u(kT)-u((k-1)T)}{T}(t-kT)
			\]
			\begin{itemize}
				\item Galat lebih kecil untuk sinyal halus.
				\item Fasa lebih buruk pada frekuensi tinggi; implementasi lebih rumit.
			\end{itemize}
			\vspace{0.1cm}
			Batas galat rekonstruksi (sinyal mulus):
			\[
				e_{\text{ZOH}} \le T\max|\dot u|,\qquad e_{\text{FOH}} \le \tfrac{T^2}{2}\max|\ddot u|
			\]
		\end{column}
		\begin{column}{0.46\textwidth}
			\centering
			\begin{tikzpicture}
				\begin{axis}[skd,width=6.6cm,height=4.6cm,xmin=0,xmax=6,ymin=-1.4,ymax=1.4,xlabel={$t/T$},ytick=\empty,legend style={at={(0.5,1.02)},anchor=south},legend columns=3]
					\addplot[skdblue!40,domain=0:6,samples=100] {sin(deg(1.1*x))};
					\addplot[skdgreen,const plot,samples at={0,1,...,6}] {sin(deg(1.1*x))};
					\addplot[skdred,samples=200,domain=1:6] {sin(deg(1.1*floor(x))) + (sin(deg(1.1*floor(x)))-sin(deg(1.1*(floor(x)-1))))*(x-floor(x))};
					\legend{asli, ZOH, FOH}
				\end{axis}
			\end{tikzpicture}\\
			{\scriptsize DAC pada praktiknya selalu berupa ZOH.}
		\end{column}
	\end{columns}
\end{frame}

% ===========================================================================
\section{Memilih Waktu Cacah\texorpdfstring{\\[4pt]{\Large\en{Choosing the Sampling Period}}}{ / Choosing the Sampling Period}}
% ===========================================================================

\begin{frame}{Pertimbangan dalam memilih $T$}
	\framesubtitle{Trade-offs in choosing $T$}
	\begin{columns}[T]
		\begin{column}{0.48\textwidth}
			\begin{alertblock}{$T$ terlalu besar}
				\small
				\begin{itemize}
					\item Tunda ZOH $T/2$ menurunkan margin fasa
					\item Respons gangguan lambat; risiko \textit{aliasing}
					\item Kinerja turun, bahkan tidak stabil (Pertemuan 1)
				\end{itemize}
			\end{alertblock}
		\end{column}
		\begin{column}{0.48\textwidth}
			\begin{alertblock}{$T$ terlalu kecil}
				\small
				\begin{itemize}
					\item Beban komputasi dan komunikasi besar
					\item Kutub diskrit menumpuk dekat $z=1$: sensitif terhadap kuantisasi koefisien (Pertemuan 15)
					\item Estimasi derivatif sangat berderau
				\end{itemize}
			\end{alertblock}
		\end{column}
	\end{columns}
	\vspace{0.3cm}
	\begin{block}{Aturan praktis \en{(rules of thumb)}}
		\small
		\begin{itemize}
			\item $\omega_s \approx 20$--$40\times$ \textit{bandwidth} kalang tertutup [FR 8.5]
			\item 4--10 sampel per waktu naik $t_r$ kalang tertutup
			\item Untuk plant FOPDT: $T \approx 0{,}1\tau$ hingga $0{,}2\tau$, dan $T$ jauh lebih kecil dari waktu tunda $\theta$
		\end{itemize}
	\end{block}
\end{frame}

\begin{frame}{Contoh: memilih $T$ untuk kontrol kecepatan motor}
	\framesubtitle{Example: choosing $T$ for motor speed control}
	\begin{columns}[T]
		\begin{column}{0.52\textwidth}
			\small
			Plant acuan semester ini: $G_v(s) = \dfrac{K}{\tau s+1}$, $K=1$, $\tau = 0{,}5$\,s.

			\vspace{0.2cm}
			Target kalang tertutup: konstanta waktu $\tau_{cl} \approx 0{,}15$\,s.
			\begin{enumerate}
				\item \textit{Bandwidth} $\omega_b \approx 1/\tau_{cl} \approx 6{,}7$ rad/s
				\item $\omega_s \approx 20\omega_b$--$40\omega_b \approx 133$--$267$ rad/s
				\item $T = 2\pi/\omega_s \approx 0{,}024$--$0{,}047$\,s
				\item Waktu naik $t_r \approx 2{,}2\tau_{cl} = 0{,}33$\,s $\Rightarrow$ 4--10 sampel: $T \approx 0{,}03$--$0{,}08$\,s
			\end{enumerate}
		\end{column}
		\begin{column}{0.44\textwidth}
			\begin{block}{Keputusan}
				\small
				$T$ di kisaran 0,03--0,05\,s wajar. Sebagai contoh analitik di kuliah dipakai $T=0{,}1$\,s agar angka mudah dihitung tangan; dampak pilihan ini dianalisis di Pertemuan 6--7.
			\end{block}
			\vspace{0.1cm}
			{\small Batasan lain: waktu eksekusi \texttt{loop()} dan pengiriman serial harus jauh di bawah $T$ (Modul 0, Percobaan 4).}
		\end{column}
	\end{columns}
\end{frame}

\begin{frame}{Ringkasan dan latihan}
	\framesubtitle{Summary and exercises}
	\begin{columns}[T]
		\begin{column}{0.48\textwidth}
			\begin{block}{Ringkasan}
				\small
				\begin{itemize}
					\item Kuantisasi: $Q=\text{FSR}/2^n$, $\sigma_q^2 = Q^2/12$.
					\item Pencacahan membuat spektrum periodik dengan periode $\omega_s$.
					\item Shannon: $\omega_s > 2\omega_m$; jika tidak, terjadi \textit{aliasing}.
					\item ZOH: $G_h(s) = (1-e^{-sT})/s$, setara tunda $T/2$.
					\item $T$ dipilih dari \textit{bandwidth} atau waktu naik kalang tertutup.
				\end{itemize}
			\end{block}
		\end{column}
		\begin{column}{0.48\textwidth}
			\begin{exampleblock}{PR-1 (ringkas)}
				\small
				\begin{enumerate}
					\item Hitung $Q$ dan $\sigma_q$ untuk ADC 12-bit $\pm 10$\,V.
					\item Sinyal 47 Hz dicacah 40 Hz. Berapa frekuensi aliasnya?
					\item Turunkan $|G_h(j\omega)|$ dan tunjukkan fasanya $-\omega T/2$.
					\item Untuk $G(s)=1/(0{,}2s+1)$ dengan target $\tau_{cl}=0{,}05$\,s, usulkan $T$ dan jelaskan alasannya.
					\item Simulasikan sinus 9 Hz dicacah 10 Hz (MATLAB/Python) dan tampilkan aliasnya.
				\end{enumerate}
			\end{exampleblock}
		\end{column}
	\end{columns}
\end{frame}

\end{document}
